(1) 在激光对称驱动下反对称态 \(|D\rangle=(|gr\rangle-|rg\rangle)/\sqrt2\) 的两条跃迁振幅相消，\(H|D\rangle=0\)，是暗态。对称基 \((|gg\rangle,|W\rangle,|rr\rangle)\) 中每一相邻跃迁耦合为 \(\hbar\Omega/\sqrt2\)，故
\[
\boxed{H_{\rm sym}=\begin{pmatrix}0&\hbar\Omega/\sqrt2&0\\
\hbar\Omega/\sqrt2&0&\hbar\Omega/\sqrt2\\0&\hbar\Omega/\sqrt2&V\end{pmatrix}}.
\]
始态完全对称，因此反对称暗态在整个理想演化中不被激发，不能将其错误计入亮态增强系数。

(2) 当 \(V/\hbar\Omega\to\infty\)，高能 \(|rr\rangle\) 动力学解耦，剩下两态有效哈密顿量为
\(H_{\rm eff}=\frac{\hbar\Omega_{\rm coll}}2(|gg\rangle\langle W|+\mathrm{h.c.})\)，
\(\boxed{\Omega_{\rm coll}=\sqrt2\Omega}\)。从 \(|gg\rangle\) 出发得
\(|\psi(t)\rangle=\cos(\Omega_{\rm coll}t/2)|gg\rangle-i\sin(\Omega_{\rm coll}t/2)|W\rangle\)，
最大转移概率1。第一个理想 \(\pi\) 脉冲长 \(\boxed{t_\pi=\pi/(\sqrt2\Omega)}\)，制得单激发Bell型 \(|W\rangle\)。

(3) 对有限但很大 \(V\)，由高能态定态振幅方程 \(V c_{rr}+(\hbar\Omega/\sqrt2)c_W\simeq0\) 得
\(c_{rr}\simeq-(\hbar\Omega/(\sqrt2V))c_W\)。回代亮态方程，产生二阶对角光移
\[
\boxed{\delta E_W\simeq-\frac{(\hbar\Omega/\sqrt2)^2}{V}
=-\frac{\hbar^2\Omega^2}{2V}},
\]
并有双激发概率 \(P_{rr}=O[(\hbar\Omega/V)^2]\)。该近似需避开瞬时开启脉冲造成的其他高频跃迁，不能把它误当所有时间的严格数值上界。

(4) 对纯 \(|W\rangle=(|gr\rangle+|rg\rangle)/\sqrt2\)，对任意一个原子求偏迹给 \(\rho_1=(|g\rangle\langle g|+|r\rangle\langle r|)/2\)。其纯度 \(1/2\)，von Neumann熵 \(S=-2(1/2)\ln(1/2)=\boxed{\ln2}\)，是最大两量子比特纯态纠缠。集体增强的 \(\sqrt2\) 来自两个单激发路径的相干叠加；理想阻塞态根本没有 \(|rr\rangle\) 占据，不是双激发光子数效应。
